\documentclass[11pt,notitlepage]{article}
\makeatletter\if@twocolumn\PassOptionsToPackage{switch}{lineno}\else\fi\makeatother

%Publisher: American Association for Cancer Research
%Template Provided By: Typeset

\usepackage{amsmath,amssymb,amstext,tabulary,graphicx,times,caption,fancyhdr,float}
\usepackage[utf8]{inputenc}
\usepackage[T1]{fontenc}
\usepackage{url,lineno,microtype,subcaption}
\usepackage[onehalfspacing]{setspace}
\usepackage{wrapfig} 
\usepackage{soul}
\usepackage{caption}
\usepackage{cite}
\usepackage{algorithmic}
\usepackage{textcomp}
\usepackage{xcolor}

\usepackage{hyperref}
\hypersetup{
    colorlinks=true,
    urlcolor=blue,
    linkcolor=black,
    filecolor=black,
    citecolor=black
}


\usepackage{abstract}
\usepackage{titlesec}
\titlespacing\section{0pt}{12pt plus 4pt minus 2pt}{0pt plus 2pt minus 2pt}
\titlespacing\subsubsection{0pt}{12pt plus 4pt minus 2pt}{0pt plus 2pt minus 2pt}
\titlespacing\subsubsubsection{0pt}{12pt plus 4pt minus 2pt}{0pt plus 2pt minus 2pt}
\usepackage[paperheight=11.66in,paperwidth=8.91in,margin=1.55cm,headsep=.5cm,top=2.5cm,columnsep=.9cm]{geometry}
\renewenvironment{onecolabstract} {\vspace*{-1pc}\trivlist\item[]\leftskip\oupIndent\hrulefill\par\vskip4pt\noindent\textbf{\abstractname}\mbox{\null}\\\hspace*{.5pc}}{\par\noindent\hrulefill\endtrivlist} 
\linespread{1.13} \date{}
\captionsetup[figure]{labelfont=sc,labelfont=normal,labelsep=period,labelfont=bf,skip=1.4pt,aboveskip=1pc}
\captionsetup[table]{labelfont=sc,labelfont=bf,skip=1.4pt,labelsep=period}
\renewcommand{\headrulewidth}{0pt}
\newcommand{\diff}[2]{\frac{d #1}{d #2}}

%%%%%%%%%%%%%%%%%%%%%%%%%%%%%%%%%%%%%%%%%%%%%%%%%%%%%%%%%%%%%%%%%%%%%%%%%%
% Following additional macros are required to function some 
% functions which are not available in the class used.
%%%%%%%%%%%%%%%%%%%%%%%%%%%%%%%%%%%%%%%%%%%%%%%%%%%%%%%%%%%%%%%%%%%%%%%%%%
\usepackage{url,multirow,morefloats,floatflt,cancel,tfrupee}
\makeatletter


\AtBeginDocument{\@ifpackageloaded{textcomp}{}{\usepackage{textcomp}}}
\makeatother
\usepackage{colortbl}
\usepackage{xcolor}
\usepackage{pifont}
\usepackage[nointegrals]{wasysym}
\urlstyle{rm}
\makeatletter

%%%For Table column width calculation.
\def\mcWidth#1{\csname TY@F#1\endcsname+\tabcolsep}

%%Hacking center and right align for table
\def\cAlignHack{\rightskip\@flushglue\leftskip\@flushglue\parindent\z@\parfillskip\z@skip}
\def\rAlignHack{\rightskip\z@skip\leftskip\@flushglue \parindent\z@\parfillskip\z@skip}

%Etal definition in references
\@ifundefined{etal}{\def\etal{\textit{et~al}}}{}


%\if@twocolumn\usepackage{dblfloatfix}\fi
\usepackage{ifxetex}
\ifxetex\else\if@twocolumn\@ifpackageloaded{stfloats}{}{\usepackage{dblfloatfix}}\fi\fi

\AtBeginDocument{
\expandafter\ifx\csname eqalign\endcsname\relax
\def\eqalign#1{\null\vcenter{\def\\{\cr}\openup\jot\m@th
  \ialign{\strut$\displaystyle{##}$\hfil&$\displaystyle{{}##}$\hfil
      \crcr#1\crcr}}\,}
\fi
}

%For fixing hardfail when unicode letters appear inside table with endfloat
\AtBeginDocument{%
  \@ifpackageloaded{endfloat}%
   {\renewcommand\efloat@iwrite[1]{\immediate\expandafter\protected@write\csname efloat@post#1\endcsname{}}}{\newif\ifefloat@tables}%
}%

\def\BreakURLText#1{\@tfor\brk@tempa:=#1\do{\brk@tempa\hskip0pt}}
\let\lt=<
\let\gt=>
\def\processVert{\ifmmode|\else\textbar\fi}
\let\processvert\processVert

\@ifundefined{subparagraph}{
\def\subparagraph{\@startsection{paragraph}{5}{2\parindent}{0ex plus 0.1ex minus 0.1ex}%
{0ex}{\normalfont\small\itshape}}%
}{}

% These are now gobbled, so won't appear in the PDF.
\newcommand\role[1]{\unskip}
\newcommand\aucollab[1]{\unskip}
  
\@ifundefined{tsGraphicsScaleX}{\gdef\tsGraphicsScaleX{1}}{}
\@ifundefined{tsGraphicsScaleY}{\gdef\tsGraphicsScaleY{.9}}{}
% To automatically resize figures to fit inside the text area
\def\checkGraphicsWidth{\ifdim\Gin@nat@width>\linewidth
	\tsGraphicsScaleX\linewidth\else\Gin@nat@width\fi}

\def\checkGraphicsHeight{\ifdim\Gin@nat@height>.9\textheight
	\tsGraphicsScaleY\textheight\else\Gin@nat@height\fi}

\def\fixFloatSize#1{}%\@ifundefined{processdelayedfloats}{\setbox0=\hbox{\includegraphics{#1}}\ifnum\wd0<\columnwidth\relax\renewenvironment{figure*}{\begin{figure}}{\end{figure}}\fi}{}}
\let\ts@includegraphics\includegraphics

\def\inlinegraphic[#1]#2{{\edef\@tempa{#1}\edef\baseline@shift{\ifx\@tempa\@empty0\else#1\fi}\edef\tempZ{\the\numexpr(\numexpr(\baseline@shift*\f@size/100))}\protect\raisebox{\tempZ pt}{\ts@includegraphics{#2}}}}

%\renewcommand{\includegraphics}[1]{\ts@includegraphics[width=\checkGraphicsWidth]{#1}}
\AtBeginDocument{\def\includegraphics{\@ifnextchar[{\ts@includegraphics}{\ts@includegraphics[width=\checkGraphicsWidth,height=\checkGraphicsHeight,keepaspectratio]}}}

\DeclareMathAlphabet{\mathpzc}{OT1}{pzc}{m}{it}

\def\URL#1#2{\@ifundefined{href}{#2}{\href{#1}{#2}}}

%%For url break
\def\UrlOrds{\do\*\do\-\do\~\do\'\do\"\do\-}%
\g@addto@macro{\UrlBreaks}{\UrlOrds}



\edef\fntEncoding{\f@encoding}
\def\EUoneEnc{EU1}
\makeatother
\def\floatpagefraction{0.8} 
\def\dblfloatpagefraction{0.8}
\def\style#1#2{#2}
\def\xxxguillemotleft{\fontencoding{T1}\selectfont\guillemotleft}
\def\xxxguillemotright{\fontencoding{T1}\selectfont\guillemotright}

\newif\ifmultipleabstract\multipleabstractfalse%
\newenvironment{typesetAbstractGroup}{}{}%

%%%%%%%%%%%%%%%%%%%%%%%%%%%%%%%%%%%%%%%%%%%%%%%%%%%%%%%%%%%%%%%%%%%%%%%%%%

\newcommand{\ILF}[0]{C_\text{IL-15}}
\newcommand{\ILS}[0]{C_\text{IL-7}}
\newcommand{\MemE}[0]{M_\text{CD8}}
\newcommand{\EffE}[0]{E_\text{CD8}}
\newcommand{\MemF}[0]{M_\text{CD4}}
\newcommand{\EffF}[0]{E_\text{CD4}}
\newcommand{\norm}[1]{\|#1\|}
\newcommand{\bignorm}[1]{\bigg\|#1\bigg\|}
\newcommand{\inner}[2]{\langle #1,#2 \rangle}
\newcommand{\avg}[1]{\langle #1 \rangle}
\newcommand{\pdiff}[2]{\frac{\partial #1}{\partial #2}}
\newcommand{\ppdiff}[2]{\frac{\partial^2 #1}{\partial #2^2}}
\newcommand{\ddiff}[2]{\frac{d^2 #1}{d #2^2}}
\newcommand{\intinf}[0]{\int_{-\infty}^{\infty}}
\newcommand{\E}[0]{\mathbb E}
\newcommand{\un}[0]{u^{(n)}}
\newcommand{\vn}[0]{v^{(n)}}
\newcommand{\En}[0]{E^{(n)}}
\newcommand{\arctanh}[0]{\text{ arctanh}}



\makeatletter\def\oupIndent{1pt}
\def\author#1{\gdef\@author{\hskip-\dimexpr(\tabcolsep)\hskip\oupIndent\parbox{\dimexpr\textwidth-\oupIndent}{\centering\bfseries#1}}}
\def\title#1{\gdef\@title{\centering\bfseries\ifx\@articleType\@empty\else\@articleType\\\fi#1}}
\let\@articleType\@empty \def\articletype#1{\gdef\@articleType{{\normalfont\itshape#1}}}

\AtBeginDocument{\fancypagestyle{plain}{\renewcommand{\headrulewidth}{0pt}\renewcommand{\footrulewidth}{0pt}\fancyhf{}}}
\makeatother



\begin{document}

\subsection*{Main}
% Polyploid/ polyaneuploid giant cancer (PGCC/PACC) (PGCC henceforth) state is a common response of cancer cells to various stressors including chemotherapy, irradiation, hypoxia and viral infection. Upon stress, PGCCs adopt an endoreplication in which the genome replicates, mitosis is omitted, and cells grow in size, leading to resistance. How accessible WGD is to a cell, therefore, directly translates to its resistance to therapy. In this study, we hypothesized that WGD and PGCC state are more plausible to cancer cells that already have a higher ploidy content prior to therapy. To test this hypothesis, we developed a comprehensive three-tier framework consisting of i) computational, ii) in vitro and iii) in silico components.
% We designed a set of ordinary differential equations (ODEs) to model how cells enter and exit the PGCC state in response to a given stressor. We engineered how this in silico approach interacts with in vitro experiments into a broadly applicable software solution called CLONEID. The software uses computer vision to monitor phenotypic changes in cell and nuclear size from standard bright-field microscopy and classify cells into PGCC and non-PGCC states. We used CLONEID to test various therapeutic agents for their ability to select for a stable near-tetraploid population in a set of near-diploid cell lines. Through spontaneous cell fusions, we also obtained tetraploid breast cancer cells matched with their parental lines. Altogether, this framework, now, enables us to monitor the PGCC state experimentally in matched cell lines with differential DNA content to model the successful entrance and exit rates to and from the endoreplication state, respectively.

\subsection*{Results}

% As the first application of this framework, we tested the ability of our matched triple-negative breast cancer (TNBC) lines (SUM-159 and MDA-MB-231) to access the PGCC state upon treatment with 18 commonly-used chemotherapy agents. We observed that gemcitabine caused continued cell growth without cell division in both tetraploid SUM-159 (Figure 1A-F) and MDA-MB-231 cells (Figure 2A) whereas near-diploid parental cells were hypersensitive to the treatment. Consequently, tetraploid cancer cells continued to safely grow in the presence of gemcitabine (Figure 3A). Furthermore, these PGCCs re-entered the proliferative cell cycle and grew in cell number when treatment is terminated (Figure 3B).

% Explain CLONEID and how it informed us that the cells turned into PACC state. 

% Explain mathematical modeling informing on the switches between polyploid state and proliferative state.

% Based on these findings, we have formulated an in silico model of the interactions of cancer cells, ploidy level and drug response as follows. Let C and P be the compartments of cycling and polyploid cells respectively. Then we can model their temporal dynamics when exposed to x units of a drug (i.e. ploidy-modulating agent) as shown in equations (1) and (2) in Figure 3.

%  Cycling and polyploid cells die at different rates (d\_C and d\_P respectively), which are both a function of drug concentration x (Figure 3). Hereby, we assume that cells cease dividing under treatment and that resistance is a function of how effectively tumor cells can enter and exit the polyploid state in response to a given drug concentration. The model can easily be expanded to also account for multiple rounds of polyploidization and for potential depolyploidization after treatment is terminated.

% Gemcitabine (GEM)-based chemotherapy is a standard treatment for patients with TNBC although its efficacy is limited mainly due to drug resistance. Moreover, such poor response rate is coupled with severe side effects that frequently leads to the failure of vital organs, serious secondary diseases, and co-morbidities. Our findings suggest that ploidy is a predictive biomarker for gemcitabine sensitivity. Thus, we expect our findings and three-component framework strategy to help stratify the TNBC patient population by their response to gemcitabine. In addition, our mathematical modelling approach has the promising potential to inform personalized dose optimization and to effectively decrease administered gemcitabine dose for a subset of patients, which would alleviate therapy-associated side effects and lethalities.


\clearpage
\subsection*{Methods}
% \subsubsection*{Model 1}

% Hypothesis 1: 
% Cells continue to divide under treatment, but most of them die in the process. Over time resistance develops as dividing cells reduce their apoptotic risk.\newline


% We use a log-kill pattern for modeling treatment effect, which assumes that the shrinkage rate of the population as a result of drug treatment is proportional to population size:


% \begin{equation} \label{m1}
% \frac{\mathrm{d} T}{\mathrm{~d} t}=f(T)-k_{\mathrm{d}} \cdot e^{-\lambda \cdot t} \cdot \text { Exposure } \cdot T
% \end{equation}

% The $e_{-\lambda \cdot t}$ term quantifies the decline of drug effect over time (i.e. resistance) and $f(T)$ is the growth term. For example, for logistic growth we have: $f(T) = r \cdot T \cdot (1 - \frac{T}{K})$

\subsubsection*{Cell viability and cell death assay}
Cell viability assays were carried out using the IncuCyte live-cell imaging system (Sartorius). Briefly, 2N and 4N SUM-159-NLS-mCherry cells were plated in four replicates in 96-well plates (Falcon, Corning, New York, USA) at a density of 1000 cells per well and treated with a 10-point 2X serial dilutions of gemcitabine within the concentration range of 0 and 800 nM. Sytox Deep Red (S11381, ThermoFisher Scientific) nucleic acid stain was also added to monitor dead cells. Plates were then incubated in the IncuCyte CO2 incubator for 10 days while four separate images were captured from each well every 2 hours using 10× objective in red, infrared and phase channels. Red object count (ROC) per well was calculated as a measure of viable cells and near-infrared object counts (NOC) per well was calculated as a measure of dead cells when the treatment was performed.

\subsubsection*{Image analysis}
 \textbf{Cell segmentation and classification with HALO}

Cell segmentation was achieved by training a single-cell segmentation neural network in Halo-AI platform [1]. At the beginning, a training dataset was collected by selecting a subsample of images from the data. Two main selection criteria were considered during the image sampling. First, the training set must have representative images from both 2N and 4N cells. Second, multiple images were sampled from different time points for each cell type 2N and 4N to ensure the capture of cell change over time during the training process. After the selection of train set, nucleus of cells was labeled using Halo built-in labeling tools. Finally, the training process was performed until all cells are accurately segmented. 
Another object classifier was trained on the top of the cell segmentation classifier results to classify each segmented cell to either dead or alive cell. At the beginning, cells are labeled either dead or alive by tracking the cells over the time and by observing changes in their nuclei. The object classifier was trained iteratively until the cell were classified correctly.

[1] HALO Image Analysis Platform version 3.6.4134 and HALO AI version 3.6.4134 (Indica Labs, Inc.), Albuquerque, NM: Indica Labs, Inc.; 2023.

\textbf{Tracking live cells with Ilastik}

Initially, all images and their classified cell masks were compiled and put together into time series and saved as tiff files. Each time series consists of a set of images/timepoints that were put in a natural temporal ordering. Subsequently, to minimize the error in cell tracking, the timepoints in each time series were sequentially co-registered by aligning each timepoint to the timepoint before it. Rigid cross correlation registration algorithm in [2] was used to align timepoints.

Lastly, Ilastik (ver 1.4.0) was used to train a structured learning cell tracking pipeline [3]. The pipeline was trained by giving the algorithm some examples of cell trajectories/lineages over time. The training examples were selected to cover all possible scenarios of cell changes such as division, dying, or merging over the course of time. The given examples were used to optimize for the weights associated with cell detections, transitions, divisions, appearances, and disappearances [4].  Moreover, other tracking parameters were adjusted to impose more restrictions on the tracking process. An example of the restrictions is applying a penalty on the appearance of new cells that weren’t a result of cell division (in some cases, cells clump together on one frame, and then they split in the next frame, then clump again the subsequent frame).  The values of the tracking parameters were selected by running the tracking several times with different values and select the parameter value that provided the best tracking results.

[2] Manuel Guizar-Sicairos, Samuel T. Thurman, and James R. Fienup, "Efficient subpixel image registration algorithms," Opt. Lett. 33, 156-158 (2008).

[3] Berg, S., Kutra, D., Kroeger, T. et al. ilastik: interactive machine learning for (bio)image analysis. Nat Methods 16, 1226–1232 (2019). https://doi.org/10.1038/s41592-019-0582-9

[4] M. Schiegg, P. Hanslovsky, B. X. Kausler, L. Hufnagel, F. A. Hamprecht. Conservation Tracking. Proceedings of the IEEE International Conference on Computer Vision (ICCV 2013), 2013.

\textbf{Track classification}

Let $x_t$ be a parental cell at the last timepoint $t$ before its cell division event gives rise to two daughter cells $a_{t+1}$ and $b_{t+1}$. We further refer to all such parental and daughter cells as pre- and post-division tracks respectively. Further we refer to all post-division tracks for which there exists a timepoint $\tau > t+1$ at which they give rise to two daughter cells, as inter-division tracks between timepoints $t+1$ and $\tau - 1$.

We identified xx inter-division tracks in untreated wells which lasted at least 18 hours (doubling time for SUM-159 is 22 hours). For XX (XX\%) of these we observed a strong correlation between time and cell area (Pearson r>=xx), suggesting these are indeed cells that progress through the cell cycle.
XX: we manually validated XX of these tracks, using the tracking accuracy measure (TRA, see http://celltrackingchallenge.net/evaluation-methodology). Validation confirmed the correctness of these tracks as well as antecedent and subsequent cell division events (accuracy = XX).\newline

We excluded low-quality tracks based on these criteria: (i) Daughter-parent cell intervals of less than 18 hours, (ii) Cell death fraction \( > 10\% \), (iii) Weak correlation between cell size and time (Pearson \( r < 0.3 \)).

\textbf{Classification of live cells to WGD and cell cycle states}

We applied a multivariate Gaussian model to track whole-genome doubling (WGD) events in cells undergoing gemcitabine treatment, assigning WGD probabilities based on cell size distributions over timepoints. Let an inter-division track $\kappa$ bedefined as \( \kappa = \{X_{i}, \dots, X_{j}\} \), representing the area of the cell at each time $i \leq \tau \leq j$.

Inter-division tracks recorded under treatment-naive conditions were used for training as follows. For each track \( \kappa \), fold changes in cell area were calculated relative to the first timepoint: $\dot{\kappa} = \{1, \frac{X_{i+1}}{X_i},\dots,\frac{ X_j}{ X_i}\}$. We used multivariate probability assessments to quantify the expected distribution of the cell area fold change of the three preceding timepoints before division:
\begin{equation}
    P(\tau = i+1 | {X_{i}, X_{i-1},X_{i-2}} ) \approx N(\mu, \sigma)
\end{equation}

We then applied the model to all post-division tracks recorded during gemcitabine exposure. A WGD event was assigned if a cell lacked division at a timepoint where division probability exceeded 0.5:
$WGD_i = WGD_{i-1} + P(\tau = i | {X_{i}, X_{i-1},X_{i-2}} )>0.5$ (see Fig. \ref{ploidyQuant}). Fold change was recalculated to the cell size assigned at the most recent cell division event.

 We calculated the proportion of WGD cells for each timepoint using the approach above, and multiplied it by the total number of live cells per timepoint to estimate the number of WGD cells at each timepoint.

We applied an SVM-based approach to classify cell cycle phases of cancer cells classified as $WGD=0$. We defined a subset of columns (features) relevant to cell morphology and nuclear characteristics to construct a cell-by-feature dataset matrix. Normalization of these features was achieved by subtracting column minimums and dividing by medians.
We implemented a Support Vector Machine (SVM) with a radial kernel to classify cells into the G1/S and G2/M cell cycle phases. Let $\kappa$ represent an inter-division track as defined above. We assume the corresponding cell remains in the G1/S phase for the first \( xx\% \) of the timepoints within \( \kappa \), transitioning into the G2M phase thereafter. This assumption provided the necessary labels for cell cycle phase classification. The trained classifier was subsequently applied to classify cell cycle phases in both post- and pre-division cell tracks.


 \begin{figure}
\includegraphics{Figs/ploidyQuant.png}
\caption{Ploidy quantification.
}
\label{ploidyQuant}
\end{figure}

\newpage

\subsubsection*{Mathematical model of polyploidization}

When exposed to Gemcitabine, cells can move between different compartments ($P_i$) via WGD ( $i \geq 0$ indicates number of WGD events). The rate at which cells undergo WGD ($k(i,x)$) or die ($a(i,x)$) depends on the drug concentration $x$ and the current ploidy state $i$. The rate at which cells proliferate also depends on drug concentration and ploidy state:

\begin{equation}
\alpha(i, x) = 
\Gamma_{\alpha} \left(1 - \frac{x}{x + u}\right)^2
\end{equation}

We assume all three rates, (cell death-, proliferation- and WGD rates), follow Michaelis-Menten kinetics as a function of drug concentration, with Michaelis constants $w, u, v$ respectively.

With this we can model the change in the concentration of live cells of different ploidy ($P_i$) and dead cells ($A$) as:


\begin{equation}\label{main}
\begin{cases}
\frac{dG}{dt} = \max(-\theta, -G) \\
\frac{dx}{dt} = -\frac{dG}{dt} - \xi \cdot a_{p,\text{tot}}(t) \cdot x \\
\frac{dP_i}{dt} = - k(i, x) P_i - a(i, x) P_i + 
k(i-1, x) P_{i-1} \\
\frac{dA}{dt} = \sum_{i=1}^{n} a(i, x) P_i + a(0, x) G1S = a_{p,\text{tot}}(t)
\end{cases}
\end{equation}

We model the cell cycle, as well as transitions from cell cycle to endocycle as:


\begin{equation}\label{cellcycle}
\begin{cases}
    \frac{dG1S}{dt} = (1 - \frac{G1S + G2M}{K}) * (2 * k_1 * G2M - k_2 * G1S) - a(0, x) G1S \\
    \frac{dG2M}{dt} = (1 - \frac{G1S + G2M}{K}) *(k_2 * G1S - k_1 * G2M) - k(0, x) * G2M \\
    \frac{dP_1}{dt} = - k(1, x) P_1 - a(1, x) P_1 + 
    k(0, x) P_{0}  \\
    \frac{dP_2}{dt} = - a(2, x) P_2 + 
    k(1, x) P_{1} 
\end{cases}
\end{equation}



% We further model the change in Gemcitabine concentration ($G$) over time as:

% \begin{subequations}\label{main_drug}
% \begin{align}
% \diff{G}{t} =-\theta \cdot \frac{G}{G + \nu} \cdot \sum_i P_i  \\
% \diff{x}{t} = \theta \cdot \frac{G}{G + \nu} \cdot \eta - \xi x
% \end{align}
% \end{subequations}

, where $G$ is the average concentration of the active component of Gemcitabine (dFdCTP) inside a cell, $\theta$ is the rate at which dFdCTP is incorporated into DNA, $x$ is dFdCTP incorporated into the cells' DNA, and $\xi$ is the rate at which $x$ leaves the system, which is proportional to the total apoptosis rate  \( a_{p,\text{tot}} \). $K$ is the carrying capacity of the environment in which the cells are grown. The total apoptosis rate is the sum of apoptosis rates normalized by the sum of the populations \( P \):

\begin{equation}
a_{p,\text{tot}} = \frac{ \frac{dA}{dt} }{\sum_{i=1}^{n} P_i}
\end{equation}

As mentioned above, both cycling and polyploid cells die at a rate dependent on the amount of the active drug present inside the cell. The drug response rate of  cells to Gemcitabine concentration is dependent on the rate at which cells convert the chemotherapy drug into the active ingredient dFdCTP. Additionally, cellular response rely on the  rate at which dFdCTP is incorporated  into the cell’s DNA. This process results in a natural delay in drug response after treatment. During  this period, concentrations levels of the drug  grow steadily intracellular. This delay is reflected in the model by an  initial steep rise in x (dFdCTP) curve seen in figure $\ref{drug_level}$  below. Upon drug  utilization,  we observe a steady decline in the amount of the active drug component until it is fully exhausted and flattens out. This pattern is observed in both 4N and 2N wells.

 \begin{figure}[!htb]
\includegraphics{Figs/dFdCTP_place_holder.png}
\caption{dFdCTP levels  (Place Holder).
}
\label{drug_level}
\end{figure}


An overview of all model parameters is given in table \ref{parameterTable}. Above system can be expanded to also model depolyploidization. \newline

\begin{table}[h!]
\centering
\begin{tabular}{rl}
  \hline
  \hline
  u & Drug concentration at which proliferation is half maximum \\
  w & Drug concentration at which apoptosis is half maximal \\
  v & Drug concentration at which WGD rate is half maximal \\
  $\Gamma_{\alpha}$ & maximum proliferation rate \\
  $\Gamma_a$ & maximum apoptosis rate \\
  $\Gamma_k$ & maximum WGD rate \\
  $\gamma_a$ & background apoptosis rate \\
  $\lambda$ & ploidy conferred decline in drug-induced apoptosis \\
  $\mu$ & ploidy conferred decline/increase in WGD rate \\
  $\theta$ & Maximum rate at which Gemcitabine is transported into cells \\
  $\nu$ & Gemcitabine concentration at which cell-transporter are at half maximal saturation \\
  $\eta$ & Rate at which Gemcitabine is converted into dFdCTP inside a cell \\
  $\xi$ & Rate at which dFdCTP is incorporated into DNA \\
   \hline
\end{tabular}
\caption{Model parameters}
\label{parameterTable}
\end{table}

\subsubsection*{Optimization routine}

We observed a higher resistance to Gemictabine-exposure among 4N- compared to 2N SUM-159 lineages (Fig. XX), further denoted as $T$ and $D$ cells respectively. Because both $T$ cells and the few $D$ cells that did survive exposure, were enlarged in response to Gemcitabine, we hypothesized that $T$ cells were more efficient at entering the polyploid state, protecting them from Gemcitabine-induced death. To test this hypothesis we fit the model given by eq. \eqref{main} to the data, assuming a monotonically decreasing relationship between the number of WGD events, $i$, and probability of cell death:

\begin{subequations}\label{death_rate}
\begin{align}
a(i, x) = \Gamma_a \left(\frac{x}{x + w}\right)^2 e^{-\lambda i}
\end{align}
\end{subequations}

, whereby the exponential term reflects the decline in drug-induced apoptosis with increasing ploidy, and $K$ is the maximum number of WGD events during the duration of the experiment. We tested various functional forms of the dependence of WGD rate on number of prior WGD events, including constant, monotonically decreasing and increasing relationships: 

\begin{subequations}\label{endo_rate}
\begin{align}
k(i, x) =  \Gamma_k \cdot \frac{{x}}{{x + v}} \\
k(i, x) =  \Gamma_k \cdot \frac{{x}}{{x + v}} \cdot e^{-\mu i} \\
k(i, x) =    \Gamma_k \cdot \frac{{x}}{{x + v}} \cdot e^{-\mu (K-i)}
\end{align}
\end{subequations}

For a given choice of the four parameters $u, w, v$, the goodness of fit is evaluated using the Wasserstein distance as a metric to measure the discrepancy between observed and simulated data. ODEs are solved using an appropriate solver and the resulting solution vectors are denoted as $D$ and $T$ for diploid and tetraploid cells respectively.

Let $\hat D, \hat T$ be the measurements corresponding to simulations $D, T$ respectively. The overall discrepancy between the simulated and observed data is then calculated as:
\[
J = \frac{1}{2K} \sum_{i=1}^{K} \text{wsd}(D_i, \hat D_i) + \text{wsd}(T_i, \hat T_i)
\]
where $\text{wsd}(Y, \hat Y)$ calculates the Wasserstein distance between the simulated and observed data. This equation represents the combined discrepancy between the simulated data for both cases (N2 and N4) and their corresponding observed data, providing a quantitative measure of the goodness of fit for the given parameter values $u, w, v$.


\subsubsection*{Modeling Results}

Simulations that best recapitulated the measurements (Fig. \ref{ODE_model}A,B), made the assumption of: (i) a monotonically decreasing relationship between the number of WGD events and probability of cell death (Fig. \ref{ODE_model}D; eq. \eqref{death_rate}) and (ii) a non-monotonic relation between number of prior WGD events and WGD rate (eq. \eqref{endo_rate}d). Hereby the rate of a subsequent WGD was maximal, when cells have skipped exactly $i:=1$ mitoses (Fig. \ref{ODE_model}E).

Allowing any of the three Michaelis-constants ($u, w, v$) to vary between the two lineages did not result in a better fit and the fitted values were similar across replicates (Fig. \ref{ODE_model}F). This suggests that 4N cells are more resistant to Gemcitabine than 2N cells because the former have already skipped a mitotic event prior to drug-exposure, rather than other intrinsic differences between the two lineages. 


 \begin{wrapfigure}{r}{0.9\textwidth}
\includegraphics[width=0.9\textwidth]{Figs/ODE_model.png}
\caption{\textbf{Tetraploid state confers resistance to Gemcitabine}. (A-B) ODE model of polyploidization (lines: $P_i$) recapitulates experimental observations (points: $M_i$) for near-diploid (A) and near-tetraploid SUM-159 lineages (B). (C-D) ODE model assumes fixed relationships between the number of WGD events (x-axes) and proliferation (C), death (D) and WGD rate (E). With the exception of proliferation rate, these relations are assumed to be identical irrespective of lineage (diploid vs. tetraploid). (F) Michaelis-constants of proliferation (a), WGD (v) and death (w) fitted across multiple replicates (n=4).
}
\label{ODE_model}
\end{wrapfigure}


\clearpage

\renewcommand{\figurename}{Supplementary Figure}
\setcounter{figure}{0}  
\renewcommand{\tablename}{Supplementary Table}
\setcounter{table}{0}  

\section*{Supplementary Information}

\subsection*{Identification of drug categories associated with ploidy-related drug sensitivity}

%% Script: GDSC_analysis.R

We integrated data from the "Genomics of Drug Sensitivity in Cancer" database (GDSC) (XX: ref) with ploidy information reported in CellPassports for XX cell lines (XX: ref). We used annotations from DrugBank (XX: ref) and PubChem (XX: ref) to assign drugs to one of six categories: (i) CYTOTOXIC, (ii) EPIGENETICS AND TRANSCRIPTION, (iii) IMMUNOSUPPRESSIVE AGENTS, (iv) SIGNALING, (v) ANTINEOPLASTIC AGENTS and (vi) HORMONES AND ANTIHORMONES. We calculated the Pearson correlation coefficient between ploidy of a cell line and drug sensitivity (IC50) of the corresponding cell line for drugs from each category (examples shown for three cancer types in Supplementary Fig. \ref{SI_GDSC_vs_ploidy}). Using the `enrichment()` function from the R package `EnrichIntersect`, we then performed enrichment analysis to identify drug categories that are significantly associated with ploidy-related drug sensitivity. \newline

For 8/24 cancers, we found a strong enrichment (p-value <=0.01) of cytotoxic agents among drugs whose IC50 was positively correlated to ploidy (i.e. low ploidy cell lines tended to be more sensitive to this class of agents than high ploidy cell lines). By contrast, a negative correlation between ploidy and IC50 of cytotoxic agents was only identified for a single cancer type (p-value <=0.01; Supplementary table \ref{SI_GDSC_vs_ploidy_A}). \newline

For 5/24 cancers we found a strong enrichment (p-value <=0.01) of agents interfering with signaling pathways among drugs whose IC50 was negatively correlated to ploidy (i.e. high ploidy cell lines tended to be more sensitive to this class of agents than low ploidy cell lines). By contrast, a positive correlation between ploidy and IC50 of agents interfering with signaling pathways was only identified for two cancer types (p-value <=0.01; Supplementary table \ref{SI_GDSC_vs_ploidy_B}).

Overall these results suggest a trend exists that links drug sensitivity to ploidy, whereby low ploidy cells tend to be more sensitive to cytotoxic agents, whereas high-ploidy cell lines have a tendency to be more sensitive to targeted therapies interfering with Signaling pathways. While this was the case for the majority of cancer type as well as in the cancer-type-agnostic analysis (Supplementary tables \ref{SI_GDSC_vs_ploidy_A}, \ref{SI_GDSC_vs_ploidy_B}), we note there are exceptions to this trend, where the relationship does not hold.


\clearpage
 \begin{figure}
\includegraphics{Figs/SI_GDSC_vs_ploidy.png}
\caption{Pearson correlation coefficient between ploidy of a given cell line and IC50 of that cell line is shown for multiple drugs administered to Breast Adenocarcinoma (BRCA), Acute Lymphoblastic Leukemia (ALL) and Ovarian cancer (OV) lines. Color code indicates drug class. 
}
\label{SI_GDSC_vs_ploidy}
\end{figure}


\clearpage
 \begin{wrapfigure}{r}{0.9\textwidth}
\includegraphics[width=0.9\textwidth]{Figs/SI_GDSC_vs_ploidy_A.png}
\caption{\textbf{Cytotoxic agents are over-represented among drugs to which low ploidy cell lines are sensitive}. We calculated the Pearson correlation coefficient between ploidy of a given cell line and IC50 of that cell line for XX drugs across XX cancer types. Shown here is the p-value at which a drug category (columns) is enriched among the 10\% of top positive correlation coefficients for a given cancer type (enrichment p-values <=0.01 are highlighted green).
}
\label{SI_GDSC_vs_ploidy_A}
\end{wrapfigure}


\clearpage
 \begin{wrapfigure}{r}{0.9\textwidth}
\includegraphics[width=0.9\textwidth]{Figs/SI_GDSC_vs_ploidy_B.png}
\caption{\textbf{Agents targeting signaling pathways are over-represented among drugs to which high ploidy cell lines are sensitive}. We calculated the Pearson correlation coefficient between ploidy of a given cell line and IC50 of that cell line for XX drugs across XX cancer types. Shown here is the p-value at which a given drug category (columns) is enriched among the 10\% of top negative correlation coefficients for a given cancer type (enrichment p-values <=0.01 are highlighted green).
}
\label{SI_GDSC_vs_ploidy_B}
\end{wrapfigure}

\clearpage

\bibliography{references_Zotero}
\bibliographystyle{nihunsrt}

\end{document}